% Folder 57, batch 06 — archivists' pages 101–120.
% Pass: Opus 5.5 (claude-opus-5-5), 2026-10-03 — first pass, unchecked against the pages by a human.
%
% Skipped: 116 and 118 (typed versos, upside down, pages "20" and "19" of a
% typescript on descent, unrelated to the notes in hand, no ink of his bearing
% on them); 103 and 114 are covers in his hand, used as section headings;
% the lower half of 113 is blank.
\documentclass[11pt,a4paper]{article}
\input{../preamble/grothendieck}

\folder{57}
\batch{6}
\pages{101}{120}
\foldertitle{Picard : tapuscrit annoté (s.d.), lettres (s.d., 1962).}
\dating{1962-[vers 1968]}
\watermark{Édition de démonstration}

\begin{document}

\page{101}
\note{la page ouvre sur un crochet : l'argument qui précède, s'il y en a un, est hors du lot}

[N.B. \emph{Lemme} $X$ préschéma, $U$ ouvert rétrocompact, $f\colon X'\to X$ morphisme
\emph{fid. plat} et \emph{qu.-cpt}, $U'=f^{-1}(X')$\note{sic : on attend $f^{-1}(U)$}, on
suppose $i'_*(\mathcal{O}_{U'})\xleftarrow{\ \sim\ }\mathcal{O}_{X'}$ (ou, \ill{},
\emph{équivalemment} : la formule $i_*(\mathcal{O}_U)\xleftarrow{\ \sim\ }\mathcal{O}_X$).
Soit $P(X)$ l'ens. des classes, à isom. près, de faisceaux loc. libres de type fini sur
$X$, définissons de même $P(X')$, $P(U)$, $P(U')$. Alors le diagramme

\begin{tikzcd}
P(X) \arrow[r] \arrow[d] & P(X') \arrow[d] \\
P(U) \arrow[r] & P(U')
\end{tikzcd}

est \emph{cartésien} \struck{\ill{}}. En particulier, ce diagramme

\begin{tikzcd}
\mathrm{Pic}(X) \arrow[r] \arrow[d] & \mathrm{Pic}(X') \arrow[d] \\
\mathrm{Pic}(U) \arrow[r] & \mathrm{Pic}(U')
\end{tikzcd}

\uncertain{l'est} \emph{\uncertain{aussi}}.

\marginal{$\underline{E}'$ ; \quad $X'\xleftarrow{\,i'\,}U'$, $f$, $g$, $X\xleftarrow{\,i\,}U$ ;
\quad $\underline{F}$}

En effet, \struck{\ill{}} le foncteur $\underline{E}\mapsto i^*(\underline{E})$ des
\uncertain{faisceaux} loc. libres sur $X$ dans les \uncertain{Modules} loc. libres sur $U$
est \emph{pleinement fidèle}, donc $P(X)\to P(U)$ est injectif. Idem de même
$P(X')\to P(U')$, d'où il \ill{} donc $P(X)\to P(U)\times_{P(U')}P(X')$ est injectif.
\struck{Reste à prouver} \add{Prouvons la} surjectivité. Soient $\underline{F}$ loc. libre
sur $U$, $\underline{E}'$ loc. libre sur $X'$, tels que
$i'^*(\underline{E}')\simeq g^*(\underline{F})$. Je dis que $\underline{E}=i_*(\underline{F})$
est loc. libre $\Longleftarrow$ \add{$f$ fid. plat qu.-cpt} $f^*i_*(\underline{F})$ est loc.
libre, Or
\[
f^*i_*(\underline{F})\simeq i'_*g^*(\underline{F})\simeq i'_*i'^*(\underline{E}')
\simeq\underline{E}' , \quad \text{OK.}
\]
car $i^*(\underline{E})\simeq\underline{F}$, [Je dis que
\[
f^*(\underline{E})\simeq\underline{E}' \Longleftarrow
i'^*f^*(\underline{E})\simeq i'^*(\underline{E}')
\]
\add{$f$ plat ; $g^*\underline{F}\simeq i'^*(\underline{E}')$ ;
$i'_*(\mathcal{O}_{U'})=\mathcal{O}_{X'}$ ; $g^*i^*(\underline{E})\simeq g^*(\underline{F})$},
o.k.

\marginal{N.B. Il suffit de se borner à des faisceaux loc. libres. On peut \uncertain{prendre}
les faisceaux $\underline{E}$ tels que prof $\underline{E}_x\geq 2$ pour $x\in X-U$, et
$\underline{E}'$ \uncertain{correspondants} \ill{} [$X$ loc. noeth.]}

\page{102}
\emph{Corollaire} Si $\mathrm{Pic}(U)\to\mathrm{Pic}(U')$ est injectif (resp. bijectif), il en
est de même de $\mathrm{Pic}(X)\to\mathrm{Pic}(X')$ ….
\note{le premier argument de $\mathrm{Pic}$ est surchargé ($U$ sur $X$, semble-t-il) ; la
lecture suit le diagramme cartésien de p.~101}

\section{Th. de Seshadri et pbs correspondants}
\note{titre pris sur la couverture, de sa main (p.~103)}

\page{104}
\emph{Théorème} (Seshadri) Soit $f\colon X\to S$ un morphisme propre, plat, \struck{\ill{}}
tel que pour tout $s\in S$, \struck{$k(s)\xrightarrow{\sim}H^0(X_s,\mathcal{O}_{X_s})$ et}
$X_s$ soit \add{(\uncertain{géom. intègre} $(k(s))$ \ill{}) $(S_2)$}. On suppose $S$ loc.
noeth. normal irréductible de pt générique $\eta$, \struck{et que $X/S$ est localement
\ill{} section, etc.}, on suppose que $\underline{\mathrm{Pic}}_{X/S}$ existe, et est
\emph{essentiellement propre} sur $S$ [ce dernier fait étant vérifié si $f$ est simple, ou
$X=S\times_k Y$, $Y$ schéma normal$/k$, $k$ un corps].

Soit $D$ un \struck{diviseur} \add{\uncertain{1-cycle}} sur $X$, tel que
a) le supp $D=Z$ ne contient aucune \struck{\ill{}} fibre de $X$, i.e. les $Z_s$ sont
\emph{rares} dans les $X_s$ ; b) $D_\eta$ est un diviseur de Cartier. Alors \emph{$D$ est un
diviseur \struck{de Cartier}}.

\marginal{N.B. Le morphisme $\mathrm{Div}^+_{X/S}\to$ \ill{} \uncertain{seulement} \ill{}
si $D\geq 0$ !}
\drawing{en marge, un petit tableau incliné, à colonnes, où se lisent $\mathrm{Div}^+$ et
quelques mots illisibles}

\ill{} \ill{} \uncertain{dire} \ill{} \ill{} $X/S$. \ill{} une \uncertain{section}. Notons
que $D_\eta$ définit un faisceau inversible $L_\eta$, donc une section \add{rationnelle} de
$\underline{\mathrm{Pic}}_{X/S}$ sur $S$, d'où un Module

\page{105}
\struck{inversible $\underline{L}$ sur $X$}. Je dis qu'il est partout défini. Soit $U$ son
ens. de définition, $T=\overline{s(U)}$, il suffit de prouver que $T\to S$ [qui est
\add{birationnel, et} surjectif d'après la propreté essentielle de
$\underline{\mathrm{Pic}}_{X/S}$] est un isom., en vertu du [\emph{Main Theorem}] il suffit
de voir que \ill{} discrètes \ill{} \ill{} \ill{}, fibres \add{\uncertain{géom.}}
\uncertain{rationnelles} sur les corps \uncertain{extérieurs} \ill{} …. Ceci résulte
du lemme suivant :

\emph{Lemme} Soit $t\in T$ au-dessus de $s\in S$, \struck{Alors} \struck{Soit} $k$ une
\struck{clôture} extension algébrique close de $k(t)$, et $\underline{L}_t$ le Module sur
$X_k$ défini par $t$. Alors $L_t$ est définissable par un diviseur de Cartier positif
$\subset Z_k$.

En effet, \ill{} $V$ un \ill{} cas d'un \uncertain{anneau} de val. discrète dominant
$\mathcal{O}_{T,t}$, $S'$ son

\page{106}
\drawing{en marge gauche, le carré de changement de base, avec $Z$ et $Z'$ au-dessus de la
ligne du haut}

\begin{tikzcd}
X \arrow[d] & X' \arrow[l] \arrow[d] \\
S & S' \arrow[l]
\end{tikzcd}

spectre, $s'$, $\eta'$ son pt fermé et générique, $X'=X\times_S S'$,
$Z'=Z\times_S S'$ [\ill{} $Z$ \ill{} \ill{} \uncertain{prof.} $\geq 2$ \ill{}
\uncertain{irréductible} \uncertain{réduit}].

\struck{\ill{}} En vertu de \emph{Chevalley} \struck{les} $Z$ est univ. ouvert sur $S$ [car
$S$ est normal, les composantes irréductibles de $Z$ dominant $S$ --- ce qui signifie au
besoin, résulte du fait que $Z$ ne contient aucune comp. irréductible d'une fibre]
\marginal{explicitez}
\struck{N.B. \add{\ill{}} les fibres de $\mathrm{Div}_{X/S}$, $Z'$ \ill{}
\ill{} \ill{} et $(S_2)$, \ill{} \ill{} \ill{}, et $X/S$ est plat et \ill{}, \ill{} la
\ill{} \ill{} les fibres.}
\note{passage barré de longs traits obliques}
\ill{} ouvert sur $S'$, donc $Z'_{s'}=(\overline{Z'_{\eta'}})_{s'}$.

Or $L_{\eta'}$ est muni d'une section \add{\uncertain{rationnelle}},
\struck{\ill{}} \add{\ill{}} diviseur de Cartier $D_{\eta'}$, \struck{c'est donc} Comme
$X'$ plat sur $S'$, les $X'_{\eta'}$ est un ouvert dense,

\page{107}
d'où aussi une section rationnelle $\varphi$ de $L'$, et en fait une
\struck{section} \add{pseudo-\struck{section}} \struck{\ill{}} \add{\uncertain{régulière}}
[car \struck{les \ill{}} \ill{} : $X'$ \uncertain{intègre}] \add{qui d'ailleurs n'est
déterminée ($L'$ étant fixé) qu'à une multiplication près par
$\Gamma(S',\mathcal{O}_{S'})^*$}. Soit $D'$ son diviseur. Alors $D'$ est un diviseur de
Cartier \add{(\uncertain{positif} : \ill{} $\Gamma(\ldots)^*$)} sur $X'$, \struck{\ill{}} je
dis qu'il est \emph{transversal \struck{aux} fibres}. En effet, comme $X'_{s'}$ est
intègre, il suffit de voir que supp $D'\not\supset X'_{s'}$. Or sinon \struck{\ill{}} $D'$
contient un \uncertain{ouvert} dense de $X'_{s'}$, \struck{\ill{}} comme $X'_{s'}$ est
intègre, $D'\cap X'_{s'}$ \struck{\ill{}} \ill{} la composante $X'_{s'}$. Il suffit donc de
remplacer $\varphi$ par $\pi^{-n}\varphi$, où $\pi$ est l'uniformisante de $V$.

\marginal{Variante. \ill{} la condition \ill{} $\varphi$ \ill{} \ill{} \ill{} $V$ \ill{}
section $\varphi$}

Donc \ill{} $L'_{s'}\simeq L(D'_{s'})$, d'autre part
$\mathrm{supp}\,D'_{s'}\subset(\overline{\mathrm{supp}\,D'_{\eta'}})_{s'}$ \struck{\ill{}}
[vrai pour tt diviseur de Cartier \emph{transversal aux fibres} \ill{}
\uncertain{Est-il} \uncertain{vrai} \ill{} \ill{} $\geq 0$, \ill{} \ill{}].
\marginal{\uncertain{Pourquoi} // \uncertain{Exact}}

\page{108}
or le dernier ens. est $(\overline{Z'_{\eta'}})_{s'}=Z'_{s'}$. Ceci prouve le lemme.

Disposant d'une section de $\underline{\mathrm{Pic}}_{X/S}/S$ et ayant une section de
$X/S$, $X/S$ étant \uncertain{cohomologiquement} \ill{}, il s'ensuit que \ill{} \ill{} un
Module inversible $\underline{L}$ sur $X$, de \ill{} \uncertain{suivant} $L_\eta$. D'autre
part on a déjà \struck{\ill{}} une section de $L_\eta$ sur $X_\eta$, de diviseur $D_\eta$,
déterminée à un \uncertain{changement} \ill{} \ill{} qu'elle coïncide avec la section
\uncertain{canonique} de $L_\eta$ \ill{} \ill{} $\mathcal{O}(D_\eta)$. On voit comme
dessus que la section \struck{\ill{}} rationnelle est une pseudo-section \add{\ill{} $D'$
son diviseur} [N.B. $X$ est intègre], \add{d'ailleurs \ill{} (la \struck{pseudo-section}
$\varphi$) coïncide} \ill{} $\sigma(S)$ \struck{\ill{}} \ill{} \ill{} la section \ill{} de
$\underline{L}'$ sur $\sigma(S^{\ldots})$, donc \struck{\ill{}} \ill{} $D'$ un \ill{}
\ill{}

\marginal{\uncertain{Précisez}, il faut \ill{} \ill{} \ill{}}

\page{109}
fibres, donc [les fibres étant intègres] $D'$ est transversal aux fibres. Enfin,
\struck{\ill{}} \struck{coïncide avec} par construction $D'_\eta=D_\eta$, donc
$\Delta=D_\Sigma-D'_\Sigma$ est un diviseur tel que $\Delta_\eta=0$, \ill{} $\Delta$ ne
contient aucune fibre. On voit facilement que \ill{} $\Delta$ est nul (\emph{Lemme} :
explicitez). On a donc $D=D'$, donc $D$ est Cartier.

\emph{Cas général} Soit $S'$ l'ouvert de $X$ formé des pts où $f$ est simple.
\struck{Il suffit} Si $X\neq\emptyset$, donc $S'\to S$ surjectif, donc on est ramené au cas
de $X'/S'$.
\marginal{\uncertain{critère} \ill{} de la base \ill{} [explicitez]}

\emph{Question 1} Supposons, au lieu de l'existence de $\underline{\mathrm{Pic}}_{X/S}$ et de
propreté essentielle, que le \emph{foncteur} $\underline{\mathrm{Pic}}_{X/S}$ ait la
propriété de \struck{\ill{}} \add{propreté essentielle} [pour la section
\uncertain{rationnelle}], p.ex. $X\to S$ \emph{simple}, \struck{\ill{} $X=S\times_k Y$,
$Y/k$ \ill{}} \struck{Est-il vrai alors} Les conclusions du Th. restent-elles vraies ?

\page{110}
\emph{Question 2} Soit $f\colon X\to S$ un morphisme \struck{\ill{}} de type fini de schémas
\uncertain{affines} noethériens, $f$ \struck{\ill{}} une intersection complète
\uncertain{relative} de dim relative $\geq 3$ \struck{p.ex.} p.ex. $f$ simple, et supposons
$S$ \add{irréductible} normal. Soit $D$ sur $X$ satisfaisant les conditions a), b) du th.
Alors est-il vrai encore que $D$ est un diviseur de Cartier ??

\emph{Corollaire du th.} Soit $R$ graphe d'équivalence \uncertain{propre} dans un
$S$-préschéma $X$, on suppose que $R\to X$ est \emph{plat} \emph{propre} : fibres
\add{géométriques} $(S_2)$ et \emph{intègres}, on suppose enfin $X$ \emph{normal}
\add{\emph{intègre}} \struck{\ill{}} [à substituer à l'hyp. $X/R$ normal, $X/R$ \ill{}
\ill{} …]. Soit $D$ un \struck{diviseur} \add{1-cycle} sur $X$,

\marginal{$R\rightrightarrows X$ au-dessus de $S$ ; et tel que $\underline{\mathrm{Pic}}_{R^{(1)}/X}$
existe et ess. propre sur $S$}
\drawing{en marge, sous un trait : $D$, $D'=D^{(1)}_R$, $D'_\eta$ au-dessus de
$X\xleftarrow{p_1}R\supset R^{(2)}_\eta$, avec $p_2\colon R\to X$ vers le bas}

\page{111}
\note{page entière barrée de trois longs traits obliques}
Or \struck{\ill{}} $D$ \add{[condition universellement satisfaite \ill{} $X$, p.ex. \ill{}
\ill{} (\ill{})]} est universellement ouvert sur $S$, en vertu du th. de Chevalley, [\ill{}]
donc si $D'$ est \ill{} \ill{} \ill{} dans $S'$, on a
$D'_{s'}\subset(\overline{D'_{\eta'}})_{s'}$. Or on a $D'_{\eta'}=\Theta'_{\eta'}$
ensemblistement, d'où $D'_{s'}=\Theta'_{s'}$ ensemblistement. OK.
\marginal{\uncertain{N'est-ce pas} que aucune composante irréductible de $D$}

Donc il existe une \struck{section} section $\bar{s}$ de $\underline{\mathrm{Pic}}_{X/S}$
sur $S$ qui prolonge $s$, et qui correspond donc à un diviseur \struck{$D$} de Cartier
positif $\Theta$ sur $X$, \uncertain{transversal} aux fibres \add{tel que
$\Theta_\eta=D_\eta$}. \struck{$\Theta_\eta=$} Considérons $D-\Theta$, il \add{cycle
divisoriel} ne contient aucune composante d'une fibre, et est nul sur $X_\eta$, donc est nul
[Lemme immédiat], d'où $D=\Theta$, cqfd.

Dans le cas général [$D$ pas néc. $\geq 0$] si $f$ est projectif, il suffit d'écrire $D$
comme \struck{\ill{}} \add{\uncertain{section}} \ill{} d'un

\page{112}
tel que \struck{$D$} a) supp $D$ ne contienne aucune fibre de $R$ sur $X$ b) le
\struck{diviseur} \add{1-cycle} induit par $D$ sur la fibre générique de $R$ sur $X$ est
\add{un diviseur} \struck{Cartier} [condition toujours satisfaite si $R\to X$ est simple
…]. Sous ces conditions, $D$ est un diviseur \struck{\ill{}} \ill{}
[\ill{} \ill{} \ill{} \ill{} $R$].

\emph{Démonstration} Soit \struck{$D^{(1)}=p_1^*(D)$} $R'$ l'ens. des $x\in R$ tels que
$p_1$ et $p_2$ soient \ill{} en $x$ (c'est un ouvert \ill{} \ill{} et simple \ill{}). On
\ill{} qu'il \uncertain{induit} \ill{} \ill{} \ill{} \ill{} \ill{} dense.
[\ill{} $R'=U_1\cap U_2$, $U_1$ \ill{} \ill{} \ill{}, \ill{} \ill{} \ill{} $p_1$ \ill{}
\ill{} \ill{} \ill{}, \ill{} \ill{} \ill{} \ill{} $U_2$ \ill{}]. On \ill{} \ill{}
\marginal{\ill{}}

\page{113}
à prouver que l'image inverse \add{\struck{\ill{}} $D^{(1)}_R=D'$} de $D$ dans $R'$ par
$p_1$, \struck{\ill{}} est Cartier. A fortiori, il suffit de le prouver pour l'image
\ill{} \ill{}.
\marginal{\ill{} $D'=D^{(1)}_R$ \ill{} $D$ \ill{} $R$ \ill{} \ill{} $p_1$ [\uncertain{existe}
\ill{}]}
Or \struck{\ill{}} on \ill{} que les hyp. a) b) \struck{\ill{}} reviennent à dire que $D'$
satisfait aux conditions du Th. relativement à $p_2\colon R'\to X$. \ill{} \ill{} $D$
\ill{} \ill{} \ill{} \uncertain{conclure}.

\section{Séparation de Picard}
\note{titre pris sur la couverture, de sa main (p.~114)}

\page{115}
\emph{Lemme 1} $A$ anneau noethérien \struck{intègre} \add{réduit}, $K$ son
\struck{corps} \add{anneau} \ill{} (\ill{} fractions), $\tilde{A}$ sa \struck{\ill{}}
\add{\uncertain{clôture}} \ill{}.

(i) [On suppose $\tilde{A}=\bigcap_{\mathfrak{p}\ \text{de rang}\ 1}A_{\mathfrak{p}}$].
\struck{\ill{} $\bar{A}$ \ill{} $A$}

Soit $f\in A$ \struck{$A$} \add{$A$ non diviseur de $0$}, $\mathfrak{p}_i$ les idéaux
premiers associés \ill{}.

\struck{(ii) \ill{} \ill{} $\mathfrak{p}_iA_{\mathfrak{p}_i}=fA_{\mathfrak{p}_i}$
(« \uncertain{multiplicités} 1 »)} \add{$fA=\bigcap\mathfrak{p}_i$}

(iii) $A/fA$ est réduit

\struck{Soit $t\in A$} Alors
\[
\boxed{A^*=A_f^*\cap\bigcap_i A^*_{\mathfrak{p}_i}}
\]
\add{dans un \ill{} \ill{} $A\to$ \ill{} \ill{} \ill{}}

\emph{Dém.} Posons $B=A_f\cap\bigcap_i A_{\mathfrak{p}_i}$. Soit $\mathfrak{p}$ premier de
rang 1, si $\mathfrak{p}\ni f$ \ill{}, $\mathfrak{p}=\mathfrak{p}_i$ pour un $i$, \ill{}
$A_{\mathfrak{p}}=A_{\mathfrak{p}_i}\supset B$ ; si $\mathfrak{p}\not\ni f$, alors
$A_{\mathfrak{p}}\supset A_f\supset B$. Donc $B\subset\tilde{A}$ \ill{} \ill{}.
\struck{D'autre part $A_f^*\cap\bigcap_i A^*_{\mathfrak{p}_i}=B^*$.} / Soit $x\in B$,
donc $x=a/f^k$, $x$ entier sur $A$. Soit $k\geq 0$ minimal. \struck{\ill{}} Soit
$x^n+c_1x^{n-1}+\cdots+c_n=0$, \ill{}
\[
a^n+f^kc_1a^{n-1}+f^{2k}c_2a^{n-2}+\cdots+f^{nk}c_n=0
\]
d'où $a^n\in f^kA$ ; si $k>0$, \ill{} \ill{} \ill{} $a^n\in fA=\bigcap\mathfrak{p}_i$,
d'où $a\in\bigcap\mathfrak{p}_i$ \struck{d'où $a\in fA$}, \add{\uncertain{condition}
\ill{} \ill{}} \ill{} en \ill{} \uncertain{cap} \ill{} \ill{} \ill{} \struck{\ill{}}
$a/f^k$
\note{l'argument se poursuit au-delà de la page ; la suite n'est pas dans le lot}

\page{117}
\struck{\emph{Corollaire} Soit $U$ un \ill{} \ill{} \ill{} $\mathrm{Spec}(A)$ \ill{}
$\mathrm{Spec}(A_f)$ et les $\mathfrak{p}_i$ ;}

\emph{Corollaire 1} Soit $X$ un schéma \emph{affine} sur $A$, \ill{} \ill{} \ill{} de
$X/A$, \ill{} aux points $\mathfrak{p}_i$. Alors \ill{} \ill{}.

\ill{} $X=\mathrm{Spec}\,A[T,T^{-1}]$, \ill{} \ill{}.

\emph{Corollaire 2} $A^*=A_f^*\cap\bigcap_i A^*_{\mathfrak{p}_i}$ \ill{} \ill{} \ill{}
\ill{} \uncertain{exacte} \ill{} \uncertain{groupes}

\note{diagramme très chargé, retracé en partie : une ligne exacte, sous laquelle des flèches
verticales et des courbes de liaison mènent à un second carré ; le dernier terme de la
première ligne est surchargé et entouré ; l'exposant noté « ? » dans le second carré est illisible}

\begin{tikzcd}[column sep=small, row sep=small, nodes={font=\scriptsize}]
0 \arrow[r] & A^* \arrow[r] & A_f^* \arrow[r] & \mathbb{Z}^I \arrow[r] & \mathbb{Z}^{(I)} \\
& K^*/A^* \arrow[r] \arrow[d] & K^*/A_f^* \arrow[d] & & \\
& \mathbb{Z} \arrow[r] & \mathbb{Z}^I & &
\end{tikzcd}

\begin{tikzcd}[column sep=small, row sep=small, nodes={font=\scriptsize}]
A_f^* \arrow[r] & \mathbb{Z}^I \\
A^*\times\mathbb{Z} \arrow[r] \arrow[u] & \mathbb{Z}^{?} \arrow[u]
\end{tikzcd}

\drawing{dans la boucle du diagramme, une tête coiffée d'un chapeau à large bord}

\emph{Corollaire 3} Si \ill{} \ill{} $f$ \ill{} \ill{}. Alors \struck{\ill{}}
$A_f^*\simeq A^*\times\mathbb{Z}$, i.e. \ill{} $t\in A_f^*$ \ill{} \ill{} de façon unique
sous la forme $\varepsilon f^n$, $\varepsilon\in A^*$, $n\in\mathbb{Z}$.
\struck{\ill{}}

\page{119}
\emph{Proposition} $X$ préschéma loc. noethérien \emph{réduit}
\marginal{(?) [il suffit $X$ réduit aux pts de $Y$] (sous la condition ($\alpha$)
\ill{})}
$f\in\Gamma(X,\mathcal{O}_X)$, $f$ non diviseur de 0 ; $Y=V(f)$.

(i) $\mathcal{O}_X/f\mathcal{O}_X$ réduit

\add{($\alpha$) On suppose} \struck{(ii)} \struck{les} $\mathcal{O}_x$ sont des « bons »
anneaux, i.e. leurs complétés \struck{\ill{}} satisfont (i) Condition de finitude (ii)
Condition des chaînes strictes.

Soient $Y_i$ les composantes irréductibles de $Y$. Soit $Z$ \emph{affine} sur $X$, et $g$ une
section \uncertain{rationnelle} de $Z$ sur $X$, telle que (i) $g$ définie dans $X_f$ (ii) $g$
définie aux pts génériques $y_i$ des $Y_i$.

\emph{Alors $g$ est définie partout.}

\marginal{contre-exemple si $Z$ \uncertain{affine} \ill{} ?? \quad $X$ \emph{réduit}
\ill{} \ill{} ??}

Ceci se \uncertain{ramène} au \ill{} \ill{} où $X=\mathrm{Spec}(A)$, $A$ local, et où
$Z=X[T]$, i.e. $g\in K$ \uncertain{anneau} des fractions de $A$.

\emph{Lemme 1} Soit $Y$ une partie fermée \ill{} de $X$ loc. noeth. réduit satisfaisant
($\alpha$), \struck{$Z$ affine sur $X$} $g$ une \struck{section} fonction rationnelle de
$X$, définie en dehors de $X-Y$ et les pts génériques des

\page{120}
composantes irréductibles de $Z$ qui sont de codim 1 dans $X$. Alors $g$ est
\uncertain{une} section \ill{} $\mathcal{O}_X$.

\emph{Lemme 2} Soit \struck{$X$ \ill{} \ill{}, $f$} \add{$X$, $f$ \ill{}} $X$ préschéma
\struck{$f$} réduit, $f$ section de $\mathcal{O}_X$ non diviseur de $0$, \add{et
\uncertain{tel} que $\mathcal{O}_X/f$ \uncertain{réduit}} $g$ une section de $\mathcal{O}_X$
sur $X_f$, entière sur $\mathcal{O}_X$ en tant que \struck{section} \uncertain{pseudo}-
\struck{\ill{}} fonction rat. sur $X$. Alors $g$ est partout définie.

\emph{Corollaire 1} Sous les conditions de la proposition, soit $g$ une fonction rationnelle
sur $X$. Pour que $g$ soit une section de $\mathcal{O}_X^*$, il faut et il suffit que
$g|X_f$ le soit, et que \ill{} $g_{y_i}\in\mathcal{O}^*_{y_i}$. \struck{\ill{}}

\emph{Corollaire 2} Diagramme exact

\begin{tikzcd}
\Gamma(X_f,\mathcal{O}^*_{X_f}) \arrow[r] & \mathbb{Z}^I \\
\Gamma(X,\mathcal{O}_X^*)\times\mathbb{Z} \arrow[u] \arrow[r] & \mathbb{Z} \arrow[u]
\end{tikzcd}

\note{fin de page et fin du lot ; le corollaire se poursuit peut-être au-delà}

\end{document}
